\chapter{The one-weak C2 family with three arbitrarily small openings}
\label{component:p81893ad08bed}
\noindent\textbf{Component }\path{p81893ad08bed}\par
{\small\noindent Source: \path{R13/baseline/baseline/baseline/baseline/baseline/advances/c2/c2_openings_and_unbounded_n.tex}\par}
{\small\noindent SHA256: \path{81893ad08bed97aa8f133eedf1f5dfa35d77c4a9500ec4add025d102f2c451fb}\par}
{\small\noindent Inventory use: CO1--CO11 complete h,b-h,r openings with gunit, full large-n normal and all tests\par}
\medskip
\section{The one-weak C2 family with three arbitrarily small openings}
\label{p81893ad08bed:sec:c2-three-openings}

This theorem removes three unit assumptions from the source of
Theorem~\ref{pec59340d644f:thm:c2-general-single-angle}.  Both positive openings $h$
and $b-h$ may approach zero, and $r$ may approach zero, making $n$
arbitrarily large.  The conclusions concern the entire original parent
and every independent test.  Source physicality remains an explicit
hypothesis; no new source classification is claimed.

\begin{theorem}[Three-opening single-angle source]
\label{p81893ad08bed:thm:c2-three-openings}
Use exactly the rational C2 source and contact of
Section~\ref{pec59340d644f:sec:c2-general-single-angle}.  Suppose $x>0$ is
infinitesimal, $k,b$ are positive bounded units,
\[
             0<r<1,\qquad 0<h<b,
\]
and $g=b-1+r+xh$ is a positive bounded unit.  Assume $n\geq0$
and physicality of the original source.  There is no lower bound on
$r,h,b-h$.  Then the entire original parent has tensor SOS at all
independent projective measurement pairs and all real or complex
projective states.  Literal $n=0$ is included.
\end{theorem}

We write $\rho=\sqrt x$, $\eta=h/b$, and retain the exact source
quantities $A_c,S_c,\Delta$ from that section.  All arguments below
take place over the independent coefficient and testing fields.
Every coefficient inverse used below is pure coefficient or pure
testing, as specified; no mixed denominator is inverted.

\paragraph{Entry and the two complementary openings.}
First the former assumptions that $1-r$ and $b-k^2r$ are units follow
from $n\geq0$.  Indeed all quantities in
\[
                 (b+r)A_cS_c\geq r\Delta
\]
are bounded, and $\Delta_0=(1+k_0^2)(b_0+1)$.  Taking residues yields
\[
 r_0\leq\frac{b_0}{b_0+k_0^2(b_0+1)}<1,
 \qquad
 b_0-k_0^2r_0\geq
 \frac{b_0^2(1+k_0^2)}{b_0+k_0^2(b_0+1)}>0.
\]
Thus $m$ is a positive bounded unit, and $D/x,\alpha/x$ are positive
bounded units, regardless of whether $r$ has zero residue.

Set
\[
 T_C=xC(b-h),\qquad T_A=xAh,\qquad H=x(C+A).
\]
The original source formulas, with their apparent opening denominators
cancelled, show that $T_C,T_A$ are bounded positive units and have the
same residue
\[
 T_0=\frac{k_0^2(b_0+r_0)}{2\Delta_0}>0.
\]
Consequently $H$ is bounded below by a positive real, although it may
be unbounded.  The four identities
\begin{align*}
 xC(1-\eta)&=T_C/b,&xA\eta&=T_A/b,\\
 xC(1-\eta)^2&=T_C(b-h)/b^2,&
 xA\eta^2&=T_Ah/b^2                                      \tag{CO1}
\end{align*}
are the reason no new scalar singularity appears when either opening
collapses.  In particular their last two terms add to a bounded unit,
with residue $T_0/b_0$.  The effective scalar coefficient is
\[
                         Q=4T_0/b_0>0.                  \tag{CO2}
\]
It is independent of the limiting value of $h$.

\paragraph{A bounded $n$: restoration with the normalized hard root.}
If $r$ is a unit, the exact formula for $n$ makes $n$ bounded.  Use the
contact graph $X_g,Y_g=\eta uW$ of the prior theorem.  All of its scalar
coefficients remain bounded: the only potentially unbounded weights
are $C,A$, and on the graph they occur solely through the last two
products in (CO1), divided by $x$ and multiplied by $u^2W^2$.
Its leading scalar is exactly the scalar in that theorem with $Q$
replaced by (CO2).  The quantities $n_0,m_0,D_*,E_*$ have their former
values.  Hence the entire six-term Laurent identity, its balanced
AM--GM decomposition, global equality characterization, and positive
contact Hessian apply without change.

Replace the old hard root $N_Y=(Y-Y_g)/\rho$ by
\[
        \widehat N_Y=\sqrt H\,N_Y=\sqrt{C+A}(Y-Y_g).
                                                               \tag{CO3}
\]
Its exact hard coefficient from the $C,A$ packets is
\[
             4D_v\frac{C+AD_u}{C+A}
       =4D_v\left(1+\frac{A}{C+A}u^2\right),             \tag{CO4}
\]
which is bounded and has residue $4D_{v,0}$ on the finite $u=0$
residue axis.  The other packet contributions to this coefficient
and to the mixed hard entry are bounded: $1/(C+A)=x/H$ and
$1/\sqrt{C+A}=\rho/\sqrt H$ are bounded and infinitesimal.

For clarity the complete first-normal verification is also retained.
The previous unnormalized $Y$ derivative has exactly the form
\[
 L_Y=P(B_2+xe_w)+B(\rho u e_z)+E e_*,\qquad
 B_2=8D_v\{(xA)\eta D_u-(xC)(1-\eta)\},
 \quad e_*=8v(D-\alpha).                                \tag{CO5}
\]
Here $P=uW/\rho$, $B=uZ$, $E=\rho Z$.  The functions $e_w,e_z$
come only from the bounded $n,m,D,\alpha$ packets and the signed
core, so remain bounded.  Equation (CO1) makes $B_2$ bounded with
identically zero residue for every finite $v$ on $u=0$; $e_*$ has
the same property.  The coefficient of $\widehat N_Y$ is exactly
$L_Y/\sqrt H$, so its bounded zero-residue lifts are preserved.
The $X$ derivative contains neither $C$ nor $A$ and is unchanged.
Thus all four-root remote certificates restore both full normals.

For matched finite-angle residues the scalar has exactly the old
form $f=p^2Q_b/x+B_b$, with
\[
 Q_b=4D_v\{(xC)(1-\eta)^2+(xA)\eta^2D_u\}.
\]
Both $Q_b$ and $B_b$ have bounded coefficients by (CO1).  The monic
divisions in $u^2-x$, $W^2-\kappa^2x$, $uW-\kappa x$, and $v-k$,
where $\kappa=b/k$, therefore give precisely the same bounded Gram
on $R,S,U,T$.  Its residue is unchanged.  The old first-normal
contact divisions use only bounded $B_2,e_w,e_z$ and exact stationary
contact; multiplication by the pure coefficient $1/\sqrt H$
preserves their zero-residue lifts.  Hence the complete six-root
Gram has strictly positive residue.

If $n_0=0$ at the double-zero test, use the previous exact graph
with $\eta_c=C/(C+A)$.  Now
\[
 xC(1-\eta_c)=xA\eta_c=\frac{xCA}{C+A},
 \qquad
 \left(\frac{xCA}{C+A}\right)_0=T_0/b_0.                \tag{CO6}
\]
These quantities are bounded units even when $H$ is unbounded.
Thus the exact $\epsilon$ reserve and five scalar roots in that
proof remain valid.  Its $Y$ normal is again divided by $\sqrt H$,
and its hard $Y$ residue is now $4$ rather than $4(C_*+A_*)$.
This proves the literal $n=0$ case as well as its infinitesimal
neighborhood without deleting a state coordinate.

\paragraph{An infinitesimal $r$: the opposite event can be unbounded.}
Suppose now $r$ is infinitesimal.  Let $N=4n+2$.  The exact source
formula gives
\[
             Nr=\frac{2(b+r)A_cS_c}{\Delta}.             \tag{CO7}
\]
Thus $Nr$ is a bounded positive unit, $N$ is unbounded, and both
$nr$ and $nr^2$ are bounded.  Use the exact-contact graph
\[
 X_g=-ruZ,\qquad Y_g=\eta uW,
 \qquad
 \widehat N_X=\sqrt N(X+ruZ),\quad
 \widehat N_Y=\sqrt{C+A}(Y-\eta uW).                    \tag{CO8}
\]
At the original contact $u=\rho,v=k,W=\rho b/k,Z=1$ this graph
has exactly the required $X,Y$ values.  The potentially large $n$
packet contributes $4nr^2D_uD_vu^2Z^2$ to the graph scalar, and
$-8nrD_uD_vuZ$ to its unscaled $X$ derivative.  Equation (CO7)
therefore controls the entire scalar and every full first normal,
not only their values at contact.

The hard residue on $(\widehat N_X,\widehat N_Y)$ is
\[
                  \operatorname{diag}(D_{v,0},4D_{v,0}). \tag{CO9}
\]
Indeed $H_{XX}/N=D_uD_v+O(1/N)$, since the remaining source
weights that enter $H_{XX}$ are $m,D$; (CO4) controls $H_{YY}$.
The mixed entry is bounded times
$1/(\sqrt N\sqrt{C+A})$ and has zero residue.

The five exact scalar monic divisions on $P,B,V,E$ from the previous
proof still apply.  Their residues, with $z=v^2$, are now
\[
 q=Q(1+z),\quad b_s=2\mu,\quad c_s=2(\mu-1)z,
 \quad d_s=4\mu v,\quad e=E_2z+E_0,                    \tag{CO10}
\]
where $\mu=2m_0+1>1$, $E_2=4(D/x)_0>0$,
$E_0=4(\alpha/x)_0>0$, and $Q$ is given by (CO2).
Define
$\mathcal D(s,z)=(qs+c_s)(b_ss+e)-d_s^2s/4$.
Exactly,
\begin{align*}
 \frac{\mathcal D(s,z)}{sz}={}&
 2Q\mu s/z+2Q\mu s
 +2(\mu-1)E_2z/s+2(\mu-1)E_0/s\\
 &+QE_0/z+QE_2z+Q(E_0+E_2)-4\mu.                    \tag{CO11}
\end{align*}
All six nonconstant Laurent coefficients are positive.  Exact
stationary contact, applied after the bounded graph substitution
$u=\rho s_1,W=\rho w_1$, gives a stationary zero at
$(s,z)=(1,k_0^2)$.  Alternatively this follows by direct substitution
of
\[
 m_0=\frac1{2(b_0-1)(b_0+1)},\quad
 Q=\frac{2k_0^2}{\Delta_0},\quad
 E_2=\frac{2b_0^2}{k_0^2\Delta_0},\quad
 E_0=\frac{2k_0^2b_0^2}{\Delta_0}.
\]
Here $b_0>1$ follows from $g_0=b_0-1>0$.
Rescale $z=k_0^2t$ and use exactly the balanced reciprocal-pair and
three-cycle identity of the previous theorem.  Its $s+1/s-2$ and
$t+1/t-2$ coefficients are strictly positive.  Thus (CO11) has a
unique zero in the positive quadrant and a positive definite Hessian
there.  This proves remote scalar strictness and the matched-contact
Schur inequality, with no appeal to numerical evidence.

For the full $X$ normal, direct differentiation and parity give
\[
 \left.F_X\right|_g=V A_X+B B_X,
\]
where the complete coefficients are
\begin{align*}
 A_X={}&-8m(1+\eta u^2)+(8D+4)(1-\eta)u^2,\\
 B_X={}&-8nrD_uD_v+8m(1-r)
       -(8D+4)v^2(1+ru^2)+4D_v(1-r).
\end{align*}
These are bounded polynomials in the bounded testing coordinates.
Consequently the coefficient of $\widehat N_X$ has bounded
zero-residue coefficients on $V,B$ because $1/\sqrt N$ is
infinitesimal.  The full $Y$ normal is (CO5), with the bounded
graph (CO8); (CO1) gives all of its required bounds and zero residues.
Its complete residual coefficients are
\begin{align*}
 e_w={}&8mv^2(1+\eta u^2)-(8D+4)(1-\eta)
                     +(8\alpha+4)D_uv^2\eta,\\
 e_z={}&v\{-8m(1-r)+(8D+4)r-(8\alpha+4)\}.
\end{align*}
This restores the remote scalar certificates, including $v_0=0$,
where (CO10) has strictly positive coefficients on $P,B,V,E$ and
zero cross residues.  No special $n_0=0$ corner occurs in this case.

At a matched residue write, exactly,
\[
 \frac{F_X|_g}{\sqrt N}=uA_1+WB_1
   =-\kappa B_1S+u(A_1+\kappa B_1).
\]
Both $A_1,B_1$ have identically zero residue for every finite $v$;
their finite monic coquotients also retain the factor $1/\sqrt N$.
Exact normal contact makes $A_1+\kappa B_1$ vanish at the source
contact.  The same four conductors $R,S,U,T$ therefore give bounded
zero-residue coefficients, exactly as in the bounded-$n$ argument.
For the $Y$ normal the previous division
$L_Y=RB_2+\rho E_2'$ has bounded $B_2,E_2'$ with identically zero
leading functions; division by $\sqrt H$ preserves those properties.
All quadratic normal coefficients are already bounded by
(CO4),(CO7),(CO9).  Thus the matched-contact Gram for the entire
parent is strictly positive in residue.  The other state charts and
the simultaneous sign copy use precisely the previous proof.

\paragraph{Every other commuting test.}
Choose the pure coefficient
\[
                         \sigma=\min(1,xC,xA).
\]
It is a positive bounded unit, since $T_C,T_A$ are units while
$h,b-h$ are bounded.  Retain the nonnegative excess packets
\[
 4(C-\sigma/x)D_v(uW-Y)^2,
 \qquad4(A-\sigma/x)D_uD_vY^2,
\]
and use the four outer core identities in the previous theorem with
$C=A=\sigma/x$.  These identities retain the entire $n$ packet;
they never use boundedness of $n$.  Their displayed hard entries
$4C_*,4A_*$ are simply replaced by $4\sigma_0$, while $m_0>0$
and $(D/x)_0>0$ are unchanged.  Every residue Gram stays strictly
positive.  This covers the finite-unit $u$ axis, both reciprocal
axes and their double corners, including literal projective ends.

The finite $u=0$ axis, those outer axes, and the inherited theorem for
noncommuting-residue tests exhaust all independent measurements.
Restore the retained packets and positive homogeneous factors, and
split complex states into real and imaginary parts.  This proves
Theorem~\ref{p81893ad08bed:thm:c2-three-openings}.  In particular the newly admitted
source scales permit either $A$ or $C$ to be arbitrarily larger than
$x^{-1}$, independently of an arbitrarily large $n$.  They cannot
both have unbounded ratios to this scale because $h+(b-h)=b$ is a unit.  There is no
prescribed comparison between these coefficient scales and any
testing infinitesimal.
